\subsection{A dependency bound conditional on the complete final input}
\label{resolution-dynamics:conditional-bound}

A random threshold map is \emph{oblivious} when it is constant on
all neighbor assignments. Put
\[
 p_i^-:=\min_\eta p_i(\eta),\qquad
 p_i^+:=\max_\eta p_i(\eta),\qquad
 \rho:=1-e^{-4\beta\Lambda}.
\]
We first need a posterior bound, stronger than the usual unconditional
probability that an update is not oblivious.

\begin{lemma}[Posterior activity bound]
\label{lem:resolution-dynamics-posterior-activity}
Conditional on any complete stationary spin trajectory, the forward
random maps are independent across update events. Each map is
nonoblivious with conditional probability at most \(\rho\).
Consequently, for any fixed set \(A\) of distinct update events and
every final configuration \(x\),
\[
 \Pr[\text{every event in }A\text{ is nonoblivious}\mid X=x]
 \le\rho^{|A|}.
 \tag{PA}
\]
\end{lemma}

\begin{proof}
Given the trajectory, the prior product map law is conditioned on one
separate constraint per event: its map sends the recorded neighbor
configuration to the recorded new spin. These constraints involve
disjoint source blocks, proving conditional independence. The initial
weight \(\mu(W^0)\) is constant under this conditioning.

At an event with recorded neighbors \(\eta\) and new spin \(+1\),
the posterior uniform variable lies in \([0,p_i(\eta)]\). Thus the
nonoblivious probability is
\[
 \frac{p_i(\eta)-p_i^-}{p_i(\eta)}
 \le1-\frac{p_i^-}{p_i^+}.
\]
For new spin \(-1\), it is at most
\(1-(1-p_i^+)/(1-p_i^-)\). These formulas are equally valid for
the finite random-map distribution, by summing its threshold intervals.

The function \(a\mapsto\log[(1+\tanh a)/2]\) has derivative in
\([0,2]\). The local fields in \(p_i\) span an interval of length
at most \(2\beta\Lambda\). The same derivative bound applies to
the logarithm of the probability of \(-1\), with the sign reversed.
Therefore
\[
 \frac{p_i^-}{p_i^+}\ge e^{-4\beta\Lambda},\qquad
 \frac{1-p_i^+}{1-p_i^-}\ge e^{-4\beta\Lambda}.
\]
This proves the uniform posterior activity bound, regardless of the
external field. Multiply it over \(A\) conditional on the trajectory,
then average over trajectories conditional on \(X=x\), obtaining
(PA). No independence conditional on \(X=x\) alone is asserted.
\end{proof}

\begin{lemma}[Conditional coupling error after whole sweeps]
\label{lem:resolution-dynamics-coupling}
For every site \(i\) and every final configuration \(x\),
\[
 \Pr[F_L(Z)_i\ne x_i\mid X=x]\le\theta^L,
 \qquad \theta=D\rho.
 \tag{CE}
\]
\end{lemma}

\begin{proof}
If the coupled final spins disagree at \(i\), their last update map
is nonoblivious, and at least one neighbor disagreed immediately before
that update. Follow such a neighbor to its most recent prior update,
and continue backwards until the initial configuration is reached.
Every visited update is nonoblivious. Update times strictly decrease,
so the events in this path are distinct.

Label each event by its sweep number and site color. Moving backwards
to a neighbor's last update decreases the sweep number by zero or one:
an earlier color was updated in the same sweep, and a later color in
the preceding sweep. A path from a final event in sweep \(L\) to
the initial boundary therefore contains at least \(L\) update events.
This observation is why whole sweeps, rather than an unspecified
global mixing time, are used.

There are at most \(D^{L-1}\) possible length-\(L\) prefixes from
the fixed final event. Their vertices and times are determined by the
schedule and the chosen neighbor sequence, independently of the maps.
By (PA) each such prefix is wholly nonoblivious with probability at
most \(\rho^L\), even conditional on all of \(X=x\). A union bound
gives \(D^{L-1}\rho^L\le\theta^L\). A prefix cannot reach the
initial boundary in fewer than \(L\) events. For an isolated site
the update map is oblivious, so the claimed bound also holds.
\end{proof}
